\chapter{Preface}
This document provides a reference manual for the library \emph{codata}.
The documentation is PDF-centric and authored in \LaTeX\ in order to provide 
a stable, high-quality printout.
The PDF edition preserves the intended layout, typography, equations, figures
and cross-references, ensuring a consistent reading experience across 
platforms and over time.
A HTML version of the documentation is rendered with \emph{make4ht}.

In addition, POSIX-style man pages are generated from the sources using 
the \emph{prep(1)} preprocessor and are provided as both terminal-native
documentation in the binaries and as a PDF document.

Additional HTML versions are also available for users who prefer quick
navigation, keyword search, or rapid access to specific sections. 
However, those versions provide only documentation of the API generated 
from the docstrings using \emph{Sphinx} and \emph{FORD}. 
They are provided as a convenient complement to the present PDF 
documentation, and should not be considered as a replacement for the 
present reference manual. 
The \emph{Sphinx} HTML version presents the APIs for Fortran, C, and Python 
--- whereas, the \emph{FORD} HTML version is Fortran-specific and is primarily
intended for Fortran users.


\noindent Related Documentation:
\begin{itemize}
\item Reference Manual [HTML]:         \url{https://milanskocic.github.io/codata}
\item Man Pages [GROFF, TXT, PDF]:     \url{https://milanskocic.github.io/codata/man/}
\item APIs with Sphinx [HTML]:         \url{https://milanskocic.github.io/codata/sphinx}
\item API with FORD    [HTML]:         \url{https://milanskocic.github.io/codata/ford}
\end{itemize}


\noindent Related resources:
\begin{itemize}
\item Binary releases:  \url{https://github.com/MilanSkocic/codata/releases}
\item \emph{prep(1)}:   \url{https://github.com/urbanjost/prep}
\item \emph{make4ht}:   \url{https://github.com/michal-h21/make4ht}
\item Sphinx:           \url{https://github.com/fortran-lang/sphinx-fortran-domain}
\item FORD:             \url{https://github.com/Fortran-FOSS-Programmers/ford}
\end{itemize}
